\structure{ОСНОВНАЯ ЧАСТЬ}

\section{Организационно-правовая часть}

\subsection{Конституция Российской Федерации}

Конституция Российской Федерации является основой правового регулирования
всех общественных отношений на территории государства и обладает высшей
юридической силой. В соответствии с частью 1 статьи 15 Конституции
Российской Федерации она имеет высшую юридическую силу, прямое действие и
применяется на всей территории Российской Федерации. Законы и иные
правовые акты, принимаемые в Российской Федерации, не должны противоречить
Конституции Российской Федерации \cite{constitution_rf}.

Для рассматриваемой темы Конституция Российской Федерации имеет значение
как базовый нормативный акт, закрепляющий фундаментальные принципы
обращения с информацией, пределы реализации прав в информационной сфере, а
также общие требования, которые должны учитываться при разработке и
применении методики верификации соответствия TLA+ спецификации и исходного
кода.

Статья 23 Конституции Российской Федерации закрепляет право каждого на
неприкосновенность частной жизни, личную и семейную тайну, защиту своей
чести и доброго имени, а также право на тайну переписки, телефонных
переговоров, почтовых, телеграфных и иных сообщений \cite{constitution_rf}.
Для темы выпускной квалификационной работы данное положение имеет
практическое значение, поскольку предлагаемая методика использует
инструментирование кода и формирование трасс исполнения. Если в процессе
трассировки, журналирования событий или агрегации данных исполнения
обрабатываются сведения, позволяющие установить содержание сообщений,
действий пользователей или иные данные ограниченного доступа, такие
действия должны осуществляться с соблюдением конституционных гарантий
неприкосновенности частной жизни и конфиденциальности сообщений.

Статья 24 Конституции Российской Федерации устанавливает, что сбор,
хранение, использование и распространение информации о частной жизни лица
без его согласия не допускаются \cite{constitution_rf}. Это положение
непосредственно связано с применением методики верификации в реальных
информационных системах. Хотя сама методика ориентирована на анализ
корректности поведения распределённой реализации и сопоставление её со
спецификацией TLA+, на практике формируемые трассы исполнения могут
содержать технические сведения, косвенно связанные с субъектами данных,
идентификаторами сеансов, сетевыми сообщениями или иными атрибутами,
позволяющими соотнести исполнение системы с конкретными лицами. В связи с
этим при проектировании и применении методики необходимо исходить из
минимизации обрабатываемых данных и недопустимости включения в трассы
избыточной информации о частной жизни.

Статья 29 Конституции Российской Федерации закрепляет право каждого
свободно искать, получать, передавать, производить и распространять
информацию любым законным способом \cite{constitution_rf}. Для данной темы
это означает, что анализ технической информации о работе программной
системы, исследование её поведения, построение формальных моделей и
проверка их согласованности с реализацией сами по себе соответствуют
конституционному принципу свободы получения и обработки информации, если
осуществляются законными способами и не нарушают установленные ограничения.
Именно в этих пределах разработка методики верификации выступает как
правомерная научно-техническая деятельность, направленная на повышение
надёжности и безопасности программного обеспечения.

Вместе с тем часть 3 статьи 17 и часть 3 статьи 55 Конституции Российской
Федерации указывают, что осуществление прав и свобод человека и гражданина
не должно нарушать права и свободы других лиц, а сами права и свободы
могут быть ограничены федеральным законом в той мере, в какой это
необходимо для защиты основ конституционного строя, нравственности,
здоровья, прав и законных интересов других лиц, обеспечения обороны страны
и безопасности государства \cite{constitution_rf}. Применительно к
рассматриваемой работе это означает, что использование средств
инструментирования, трассировки и анализа поведения программной системы не
может осуществляться произвольно. Если программная система функционирует в
среде, где обрабатывается конфиденциальная информация, персональные данные
или иные сведения ограниченного доступа, методика верификации должна быть
организована таким образом, чтобы не создавать угрозу нарушения прав
субъектов данных и не противоречить специальным требованиям федерального
законодательства.

Таким образом, Конституция Российской Федерации задаёт базовые правовые
ориентиры для разработки и применения методики верификации соответствия
TLA+ спецификации и исходного кода. С одной стороны, она допускает
осуществление научной, технической и аналитической деятельности по
исследованию поведения программных систем, а с другой стороны --- требует
соблюдения прав граждан на неприкосновенность частной жизни,
конфиденциальность сообщений и законные ограничения при работе с
информацией. Следовательно, разрабатываемая методика должна применяться с
учётом конституционных принципов законности, соразмерности и защиты прав
человека \cite{constitution_rf}.

\subsection{Федеральный закон от 27.07.2006 №149-ФЗ <<Об информации, информационных технологиях и о защите информации>>}

Федеральный закон №149-ФЗ <<Об информации, информационных технологиях и о
защите информации>> является одним из основных нормативных правовых актов,
регулирующих отношения в сфере информации, информационных технологий и её
защиты. Для темы данной выпускной квалификационной работы этот закон имеет
особое значение, поскольку предлагаемая методика верификации соответствия
TLA+ спецификации и исходного кода связана со сбором, обработкой,
хранением и анализом информации о поведении программной системы
\cite{fz149_info}.

Согласно статье 1 Федерального закона №149-ФЗ, закон регулирует
отношения, возникающие при осуществлении права на поиск, получение,
передачу, производство и распространение информации, при применении
информационных технологий, а также при обеспечении защиты информации
\cite{fz149_info}. Это позволяет отнести разрабатываемую методику к сфере
действия данного закона, поскольку в её рамках используются программные
средства трассировки, агрегирования данных исполнения и модельной
проверки.

Статья 3 Федерального закона №149-ФЗ закрепляет основные принципы
правового регулирования в данной области. К ним относятся законность
обращения с информацией, достоверность информации, своевременность её
предоставления, а также необходимость защиты информации и соблюдения прав
граждан при работе с ней \cite{fz149_info}. Применительно к теме работы
это означает, что сбор трасс исполнения и их последующий анализ должны
осуществляться только в рамках законных целей, а полученные результаты
должны быть достоверными и защищёнными от искажения.

Важное значение имеет также статья 5 Федерального закона №149-ФЗ, в
соответствии с которой информация рассматривается как объект правовых
отношений \cite{fz149_info}. Для рассматриваемой методики это означает,
что исходный код, формальная спецификация, локальные и глобальные трассы
исполнения, а также результаты проверки средствами TLC могут
рассматриваться как самостоятельные информационные объекты, в отношении
которых должен соблюдаться установленный режим доступа, хранения и
использования.

Наиболее значимой для практического применения методики является статья 16
Федерального закона №149-ФЗ, посвящённая защите информации. В ней
установлено, что защита информации включает правовые, организационные и
технические меры, направленные на предотвращение неправомерного доступа,
уничтожения, модифицирования, блокирования, копирования и распространения
информации \cite{fz149_info}. Для данной работы это особенно важно,
поскольку искажение трасс исполнения, их неполнота или несанкционированное
изменение могут привести к неверным выводам о соответствии реализации
формальной спецификации. Следовательно, при использовании предлагаемой
методики необходимо обеспечивать целостность, достоверность и
защищённость всех данных, формируемых в процессе верификации.

Кроме того, статья 17 Федерального закона №149-ФЗ закрепляет, что
нарушение требований законодательства в области информации,
информационных технологий и защиты информации влечёт ответственность в
соответствии с законодательством Российской Федерации \cite{fz149_info}.
Это положение подчёркивает, что при практическом применении методики
необходимо учитывать не только технические аспекты трассировки и проверки,
но и правомерность обращения с формируемыми данными.

Таким образом, Федеральный закон №149-ФЗ формирует общую правовую основу
для разработки и применения методики верификации соответствия TLA+
спецификации и исходного кода. Закон определяет сферу регулирования
отношений, связанных с информацией и информационными технологиями,
рассматривает информацию как объект правовых отношений и устанавливает
требования к её защите. Следовательно, при разработке и применении
предлагаемой методики необходимо обеспечивать законность обработки данных
исполнения, соблюдение режима доступа к информации и защиту результатов
верификации от несанкционированного воздействия \cite{fz149_info}.

\subsection{Указ Президента Российской Федерации от 05.12.2016 №646 <<Об утверждении Доктрины информационной безопасности Российской Федерации>>}

Указ Президента Российской Федерации от 05.12.2016 №646 утвердил
Доктрину информационной безопасности Российской Федерации, которая
представляет собой систему официальных взглядов на обеспечение
национальной безопасности Российской Федерации в информационной сфере
\cite{doctrine_ib_rf}. Для темы данной выпускной квалификационной работы
этот документ представляет интерес как нормативная основа, определяющая
значение информационной безопасности, надёжности функционирования
информационных систем и развития безопасных отечественных информационных
технологий.

Доктрина информационной безопасности Российской Федерации рассматривает
информационную безопасность как важный элемент национальной безопасности и
подчёркивает необходимость устойчивого и бесперебойного функционирования
информационной инфраструктуры \cite{doctrine_ib_rf}. Для рассматриваемой
темы это имеет особое значение, поскольку методика верификации
соответствия TLA+ спецификации и исходного кода направлена на повышение
корректности поведения программных систем, выявление расхождений между
формальной моделью и реализацией и, как следствие, на снижение риска
ошибок в критически важных сценариях работы распределённых систем.

Существенным положением Доктрины является акцент на развитии и
совершенствовании информационных технологий, а также на необходимости
повышения защищённости и качества программных средств, используемых в
информационной инфраструктуре \cite{doctrine_ib_rf}. В этом контексте
разрабатываемая методика может рассматриваться как инструмент,
способствующий достижению указанных целей, поскольку она позволяет
подтверждать допустимость наблюдаемого поведения реализации с точки зрения
формальной спецификации и своевременно обнаруживать противоречия между
моделью и программным кодом.

Кроме того, в Доктрине отмечается значимость обеспечения защиты прав и
свобод человека и гражданина в информационной сфере, а также важность
применения организационных и технических мер, направленных на повышение
уровня информационной безопасности \cite{doctrine_ib_rf}. Для данной
работы это означает, что методика верификации должна рассматриваться не
только как средство повышения качества программного обеспечения, но и как
часть более общего подхода к обеспечению надёжности и безопасности
информационных систем.

Таким образом, Доктрина информационной безопасности Российской Федерации
формирует стратегический контекст для разработки и применения методики
верификации соответствия TLA+ спецификации и исходного кода. Она
подчёркивает значимость обеспечения устойчивого функционирования
информационных систем, развития безопасных информационных технологий и
использования современных методов, направленных на повышение надёжности и
защищённости программного обеспечения \cite{doctrine_ib_rf}.

\subsection{ГОСТ Р 56939--2024 <<Защита информации. Разработка безопасного программного обеспечения. Общие требования>>}

ГОСТ Р 56939--2024 <<Защита информации. Разработка безопасного
программного обеспечения. Общие требования>> является важным нормативным
документом для организационно-правового анализа темы данной выпускной
квалификационной работы. В отличие от общих нормативных правовых актов,
регулирующих сферу информации и её защиты, данный стандарт непосредственно
связан с процессом создания безопасного программного обеспечения и
устанавливает требования к организации работ на различных этапах его
жизненного цикла \cite{gost_56939_2024}.

Для темы разработки методики верификации соответствия TLA+ спецификации и
исходного кода данный стандарт представляет особый интерес, поскольку
рассматриваемая методика направлена на подтверждение корректности
реализации, выявление расхождений между формальной моделью и программным
кодом, а также на повышение надёжности программной системы. Использование
формальных спецификаций, трасс исполнения и средств модельной проверки
соответствует общей задаче безопасной разработки, заключающейся в
предупреждении ошибок, способных привести к нарушениям функционирования,
уязвимостям и иным недостаткам программного обеспечения.

ГОСТ Р 56939--2024 устанавливает общие требования к содержанию и порядку
выполнения работ, связанных с созданием безопасного программного
обеспечения и устранением выявленных недостатков, в том числе
уязвимостей \cite{gost_56939_2024}. Применительно к рассматриваемой работе
это означает, что методика верификации может рассматриваться как один из
инструментов, поддерживающих выполнение требований безопасной разработки,
поскольку она позволяет анализировать соответствие поведения реализации
заданной спецификации и своевременно выявлять противоречия между моделью и
исходным кодом.

Практическая значимость данного стандарта для темы выпускной
квалификационной работы заключается также в том, что он ориентирует
разработчика на системный подход к обеспечению безопасности программного
обеспечения. Верификация соответствия реализации формальной спецификации в
этом контексте выступает не как изолированная исследовательская процедура,
а как часть более общего процесса обеспечения качества, надёжности и
безопасности разрабатываемого программного продукта.

Таким образом, ГОСТ Р 56939--2024 формирует нормативную основу для
рассмотрения предлагающейся методики как элемента безопасной разработки
программного обеспечения. Его положения подтверждают, что применение
средств формального анализа, проверки поведения системы и выявления
расхождений между моделью и реализацией соответствует современным
подходам к созданию защищённого и надёжного программного обеспечения
\cite{gost_56939_2024}.

\subsection{Приказ ФСТЭК России от 01.12.2023 №240 <<Об утверждении Порядка проведения сертификации процессов безопасной разработки программного обеспечения средств защиты информации>>}

Приказ ФСТЭК России от 01.12.2023 №240 утверждает Порядок проведения
сертификации процессов безопасной разработки программного обеспечения
средств защиты информации \cite{fstek_240_2023}. Для темы данной выпускной
квалификационной работы этот документ представляет особый интерес,
поскольку он связывает требования безопасной разработки программного
обеспечения с процедурами их официального подтверждения.

В контексте рассматриваемой темы значение данного приказа состоит в том,
что предлагаемая методика верификации соответствия TLA+ спецификации и
исходного кода может рассматриваться как один из практических инструментов,
используемых в процессе безопасной разработки программного обеспечения.
Методика направлена на выявление расхождений между формальной моделью и
реализацией, а значит способствует повышению корректности программного
обеспечения и снижению вероятности ошибок, способных повлиять на его
безопасность и надёжность.

С практической точки зрения приказ ФСТЭК России №240 показывает, что
процессы безопасной разработки подлежат формализации, оценке и
подтверждению на соответствие установленным требованиям
\cite{fstek_240_2023}. Это особенно важно для данной работы, поскольку
разработка методики верификации не является изолированной исследовательской
задачей, а может быть включена в более широкий процесс обеспечения
качества и безопасности программного обеспечения.

Таким образом, приказ ФСТЭК России №240 дополняет положения
ГОСТ Р 56939--2024 и позволяет рассматривать предлагаемую методику
верификации как часть современного подхода к безопасной разработке
программного обеспечения. Применение методов формального анализа,
трассировки исполнения и проверки соответствия реализации спецификации
согласуется с направлением развития процессов безопасной разработки,
которые в ряде случаев подлежат официальной сертификации
\cite{fstek_240_2023}.
